\documentclass[10pt,a4paper]{amsart}
\usepackage[margin=19mm]{geometry}
\usepackage{amsmath,amssymb,mathtools}
\usepackage[scr=rsfs,cal=euler]{mathalfa}
\usepackage{xfrac}
\usepackage[unicode,hidelinks,bookmarks=false]{hyperref}
\newcommand{\ie}{i.e.\ }
\newcommand{\cf}{cf.\ }
\newcommand{\resp}{respectively\ }
\newcommand{\Subsection}{\subsection}
\providecommand{\nobreakdash}{\nobreak\hbox{-}}
\setlength{\emergencystretch}{2em}
\multlinegap0pt
\newcommand{\CT}{\mathcal T}
\theoremstyle{plain}
\newtheorem{thm}{Theorem}[section]
\newtheorem{cor}[thm]{Corollary}
\newtheorem{lem}[thm]{Lemma}
\newtheorem{prop}[thm]{Proposition}
\newtheorem{conditional}[thm]{Conditional theorem}
\newtheorem{maintheorem}{Theorem}
\theoremstyle{definition}
\newtheorem{defn}[thm]{Definition}
\newtheorem{example}[thm]{Example}
\newtheorem{proposedthm}[thm]{Proposed theorem}
\newtheorem{proposedprop}[thm]{Proposed proposition}
\newtheorem{construction}[thm]{Construction problem}
\theoremstyle{remark}
\newtheorem{rem}[thm]{Remark}
\newtheorem{question}[thm]{Question}
\title[Derived Smooth and Banach Higher Groupoids: Representability and Descent]{Derived Smooth and Banach Higher Groupoids: Representability and Descent}
\author{Qingyun Zeng}
\date{Source: arXiv:2609.28220v1 (CC BY 4.0)}
\begin{document}
\maketitle

\noindent\emph{Source and licence.} The delivered work is the article shown in the title above, version arXiv:2609.28220v1 (CC BY 4.0), distributed under the Creative Commons Attribution 4.0 International licence, as recorded in the arXiv licence metadata for this exact version and shown on the abstract page saved with this item. The full licence notice, which carries the licence URL and the address of that abstract page, is the bundled file \texttt{licenses/arxiv-2609.28220-CC-BY-4.0.txt}.
\par\smallskip

\noindent\textit{Editorial scope.} Counted unit: the article's Lemma (source label lem:enriched-diagonal-local-lifting) (source0.tex lines 4878--4884, source label \texttt{lem:enriched-diagonal-local-lifting}): ``Enriched finite local anodyne lifting Under the Reedy and local-component horn hypotheses of Theorem , has local strict lifts against every finite anodyne inclusion.''. It is delivered together with the article's complete original proof (source0.tex lines 4886--4945, one proof environment adjacent to the statement, 60 lines), copied line for line from the article's own text. Required context retained uncounted: lem:derived-hypercomplete-base-change (lines 2296--2329); lem:enriched-relative-matching (lines 4751--4765). Editorial formatting only: an AMS \texttt{amsart} preamble that restates the article's own macro definitions and its own theorem declarations, and the theorem environments the retained lines use; the article's own class file is neither shipped nor input; the retained environments are renumbered sequentially in order of appearance, whereas the article prints them with its own numbering; the article's equation numbering is not reproduced and every cross-reference inside the retained text resolves to the wrapper's numbering. No citation occurs inside any retained line, so the wrapper carries no bibliography; All retained bodies are exact copies of the bundled \texttt{sources/211564/source0.tex}, so no omission receipt applies. No asset-inclusion command occurs inside any retained span and the package ships no image asset.


\section*{Retained context: lem:derived-hypercomplete-base-change (source0.tex lines 2296--2329)}

\begin{thm}[Enriched realization base change]
Let $\mathcal T$ be the small Kan-enriched site under
consideration, and let $a_{\mathrm{hyp}}$ denote its
hypersheaf localization.
\label{lem:derived-hypercomplete-base-change}
A \emph{local surjection} means
a map inducing an epimorphism on the sheaf of connected
components, equivalently an effective epimorphism in the
associated infinity-topos.

Use the global injective enriched-presheaf model or its
ordinary-sheaf or hypersheaf localization. Let $X,Y,Z$
be Reedy-fibrant strict outer diagrams, let $X\to Y$ be
a Reedy model fibration whose canonical relative homotopy
horn maps are local surjections in every positive outer
degree, and let $Z\to Y$ be any map. Then the canonical
comparison, natural in this cospan,
\[
 a_{\mathrm{hyp}}\operatorname{diag}(X\times_Y Z)
 \longrightarrow
 a_{\mathrm{hyp}}\operatorname{diag}X
 \times_{a_{\mathrm{hyp}}\operatorname{diag}Y}
 a_{\mathrm{hyp}}\operatorname{diag}Z
\]
is an equivalence. The pullback on the left is the ordinary
Reedy-model pullback and therefore computes the levelwise
homotopy pullback. Outer amplitude is unrestricted.
The geometric lifting hypotheses concern positive relative
horns; the ambient Reedy-fibration condition applies in
every degree, including zero. Absolute horns of the
diagrams are unrestricted, and the site may lack enough
points. Write $\mathbf{PB}_\tau(\mathcal T)$ for this
base-change property.
\end{thm}


\section*{Retained context: lem:enriched-relative-matching (source0.tex lines 4751--4765)}

\begin{lem}[Relative matching and strict local lifts]
Let $f:X\to Y$ be a Reedy model fibration between
Reedy-fibrant strict outer diagrams. For a finite inclusion
$A\hookrightarrow B$, the map
\[
 \{B,X\}\longrightarrow
 \{A,X\}\times_{\{A,Y\}}\{B,Y\}
\]
is a model fibration and its target computes the indicated
homotopy relative target.
\label{lem:enriched-relative-matching}
If a positive relative homotopy horn of $f$ is locally
surjective on components, every inner simplex of its
actual relative horn target has local strict lifts.
\end{lem}


\section{The counted unit: Lemma (source label lem:enriched-diagonal-local-lifting) and its complete original proof}

\begin{lem}[Enriched finite local anodyne lifting]
Under the Reedy and local-component horn hypotheses of
Theorem~\ref{lem:derived-hypercomplete-base-change},
$\operatorname{diag}f$ has local strict lifts against
every finite anodyne inclusion.
\label{lem:enriched-diagonal-local-lifting}
\end{lem}

\begin{proof}
For a single diagonal horn, use the factorization
$L\subset A\subset D$ just proved.
Lemma~\ref{lem:enriched-relative-matching} lifts the
first horizontal horn locally. The second step lifts
against the global Reedy fibration. This constructs
the filler directly, before testing Kan-ness of
either diagonal.

We must retain homotopies of test arrows when composing
covering sieves. The functor $\delta_!$ preserves
monomorphisms: it takes a boundary inclusion to the
union of the corresponding external-square faces, and
every monomorphism is built by boundary attachments.
Consequently $h_s\odot\delta_!K\to h_s\odot\delta_!L$
is a Reedy cofibration for a finite inclusion $K\to L$.
A path between test arrows $v_0,v_1:s\to t$ induces
an inner homotopy of their entire restricted squares.
If a lift at $v_0$ is given, extend it across
\[
 \bigl((h_s\odot\delta_!L)\otimes\{0\}\bigr)
 \mathop{\cup}_{(h_s\odot\delta_!K)\otimes\{0\}}
 \bigl((h_s\odot\delta_!K)\otimes\Delta[1]\bigr)
 \longrightarrow (h_s\odot\delta_!L)\otimes\Delta[1].
\]
The last tensor is inner. This is the pushout product
of a Reedy cofibration with an endpoint trivial
cofibration, hence is Reedy trivial. Lifting against
$f$ transports the filler to $v_1$. The other endpoint
and finite chains of paths give invariance under the
test-arrow class. Together with strict precomposition,
this makes \emph{existence} of a lift a sieve on
$\operatorname{Ho}(\CT)$. Only existence is required
locally; compatibility of the chosen fillers is
unnecessary.

For a finite composite of horn attachments, choose
successive local fillers. The composites of their
covering sieves generate a covering sieve by transitivity.
Homotopy invariance handles actual representatives of
composite classes that are not the selected strict
composites.

Finally a finite anodyne inclusion $K\hookrightarrow L$
is a retract of a finite relative horn-cell complex.
To see the finite support, factor it by the countable
horn small-object construction as
$K\to W\to L$, with the second map a Kan fibration.
Its lifting property gives a section $s:L\to W$
relative to $K$. Finiteness of $L$ puts $s(L)$ in
one finite stage and finitely many cells. Include the
finitely many predecessor cells supporting their finite
attaching horns. Recursing strictly decreases the
stage number, so it terminates with a finite relative
horn-cell subcomplex containing $s(L)$. Restrict the
retraction to that subcomplex. Lifting there proves
the assertion for $K\to L$.
Each problem uses only finitely many refinements,
and its covering sieve may depend on that problem.
\end{proof}

\end{document}
